\chapter*{ЗАДАНИЕ}
\addcontentsline{toc}{chapter}{ЗАДАНИЕ}

\section*{Цель работы}

Исследовать уязвимости шифра Виженера. Разработать программу для автоматического криптоанализа с использованием методов Казиски и Кирхгофа. Провести оценку криптостойкости в зависимости от длины ключа.

\section*{Задания}

\textbf{Задание 1 (Подготовка и шифрование):}
\begin{enumerate}
    \item Реализовать шифр Виженера с поддержкой русского и английского алфавитов.
    \item Зашифровать текст объемом не менее 100000 символов с использованием короткого (3-5 символов) и длинного (15+ символов) ключа.
    \item Провести визуальный и статистический анализ результатов.
\end{enumerate}

\textbf{Задание 2 (Криптоанализ --- определение длины ключа):}
\begin{enumerate}
    \item Получить у преподавателя криптограмму, зашифрованную шифром Виженера.
    \item Реализовать алгоритм Казиски для поиска повторяющихся n-грамм (последовательностей из 3+ символов).
    \item Вычислить наибольший общий делитель (НОД) расстояний между повторениями для гипотезы о длине ключа.
    \item Применить тест Куллабака-Лейблера или индекс совпадений для проверки гипотез о длине ключа.
\end{enumerate}

\textbf{Задание 3 (Криптоанализ --- взлом ключа):}
\begin{enumerate}
    \item Используя найденную длину ключа L, разбить криптограмму на L групп, каждая из которых зашифрована шифром Цезаря.
    \item Для каждой группы провести частотный анализ, сопоставляя частоты букв с эталонными для языка.
    \item Автоматизировать подбор сдвига для каждой позиции ключа, используя метод хи-квадрат для оценки совпадения с частотным профилем.
    \item Восстановить ключ и исходный текст.
\end{enumerate}

\section*{Контрольные вопросы}

\begin{enumerate}
    \item Почему метод Казиски не всегда точен и как можно повысить его надежность?
    \item Как длина ключа влияет на сложность взлома шифра Виженера?
    \item Какие современные шифры являются наследниками идеи полиалфавитности?
\end{enumerate}
